\documentclass[11pt]{article}
\usepackage[a4paper,margin=20mm]{geometry}
\usepackage{amsmath,amssymb,xcolor,graphicx}
\usepackage[hypcap=false]{caption}
\usepackage[hidelinks,bookmarksnumbered,bookmarksopen]{hyperref}
\hypersetup{pdftitle={Chemical kinetics — Rate laws and the Arrhenius equation},pdfauthor={Lecturer: Dr M. Haddad},
  pdfcreator={KherveNote}}
\renewcommand\labelitemii{\textbullet}
\renewcommand\labelitemiii{\textbullet}
\renewcommand\labelitemiv{\textbullet}
\renewcommand\labelenumii{\arabic{enumi}.\arabic{enumii}.}
\renewcommand\labelenumiii{\arabic{enumi}.\arabic{enumii}.\arabic{enumiii}.}
\renewcommand\labelenumiv{\arabic{enumi}.\arabic{enumii}.\arabic{enumiii}.\arabic{enumiv}.}
\setlength{\parindent}{0pt}
\setlength{\parskip}{0.6em}
\definecolor{knoteaccent}{HTML}{1A6DD8}
\definecolor{knotemuted}{HTML}{6B6B6B}
\definecolor{knotekey}{HTML}{D96B00}
\newcommand\knotetime[1]{\leavevmode\llap{\color{knotemuted}\scriptsize #1\hspace{1em}}}
\newcommand\knotekey[2][]{\par\noindent#1\colorbox{knotekey!12}{\parbox{\dimexpr\linewidth-2\fboxsep}{%
  \textbf{\color{knotekey}$\star$ Key point.} #2}}\par}
\newcommand\knotequestion[2][]{\par\noindent#1{\color{knoteaccent}\textbf{?}}~\emph{#2}\par}
\newenvironment{knotetranscript}{\par\color{knotemuted}\small}{\par}
\newcommand\knotesummary[1]{\par\noindent\fcolorbox{knoteaccent}{knoteaccent!6}{%
  \parbox{\dimexpr\linewidth-2\fboxsep-2\fboxrule}{\textbf{Summary}\par #1}}\par}
\newenvironment{knotechunk}{}{}
\pagestyle{plain}
\begin{document}
\begin{knotechunk}
{\color{knoteaccent}\rule{\linewidth}{1.5pt}}\par
{\LARGE\bfseries Chemical kinetics — Rate laws and the Arrhenius equation}\par
{\large Lecturer: Dr M. Haddad}\hfill{\color{knotemuted}08 October 2026 \textperiodcentered{} Chemistry building, LT1}\par
{\color{knoteaccent}\rule{\linewidth}{0.6pt}}\par
\knotesummary{Rate laws are found by experiment, not from the balanced equation. Integrated rate laws give straight-line plots that reveal the order; the Arrhenius equation links the rate constant to the activation energy, and a catalyst works by offering a path with a lower one.}

\section{Rate laws}

\knotetime{01:25}For $aA + bB \rightarrow$ products, the rate is $r = -\frac{1}{a}\frac{d[A]}{dt} = k[A]^m[B]^n$.

\knotekey[\knotetime{01:50}]{The orders $m$ and $n$ are measured; they need not equal $a$ and $b$.}

\begin{itemize}\setlength{\itemsep}{0pt}\setlength{\parskip}{0pt}
  \item overall order $m + n$; units of $k$ depend on it
  \item method of initial rates: change one concentration, watch the initial rate
\end{itemize}
\end{knotechunk}

\begin{knotechunk}
\section{Integrated rate laws}

\begin{enumerate}\setlength{\itemsep}{0pt}\setlength{\parskip}{0pt}
  \item zero order: $[A] = [A]_0 - kt$; $t_{1/2} = [A]_0/(2k)$
  \item first order: $\ln[A] = \ln[A]_0 - kt$; $t_{1/2} = \ln 2 / k$, independent of $[A]_0$
  \item second order: $\dfrac{1}{[A]} = \dfrac{1}{[A]_0} + kt$; $t_{1/2} = 1/(k[A]_0)$
\end{enumerate}
\begin{itemize}\setlength{\itemsep}{0pt}\setlength{\parskip}{0pt}
  \item plot $[A]$, $\ln[A]$ and $1/[A]$ against $t$: the straight one gives the order
  \item radioactive decay and many isomerisations are first order
\end{itemize}
\end{knotechunk}

\begin{knotechunk}
\section{Temperature: the Arrhenius equation}

\knotetime{24:25}\[k = A\exp\left(-\frac{E_a}{RT}\right), \qquad \ln k = \ln A - \frac{E_a}{R}\cdot\frac{1}{T}\] A plot of $\ln k$ against $1/T$ has slope $-E_a/R$.

\subsection{Worked example}

\knotetime{25:15}With $E_a = 50$ kJ mol$^{-1}$, going from 298 K to 308 K: $\ln(k_2/k_1) = \frac{E_a}{R}\left(\frac{1}{T_1} - \frac{1}{T_2}\right) = \frac{50000}{8.314}\times 1.09\times10^{-4} \approx 0.66$, so $k_2/k_1 \approx 1.9$.

\knotekey[\knotetime{25:40}]{The "rate doubles every 10 °C" rule of thumb holds for $E_a$ around 50 kJ mol$^{-1}$ near room temperature.}
\end{knotechunk}

\begin{knotechunk}
\section{Mechanisms}

\begin{itemize}\setlength{\itemsep}{0pt}\setlength{\parskip}{0pt}
  \item elementary steps have orders equal to their molecularity
  \item the rate-determining step sets the rate law
  \item steady-state approximation for reactive intermediates: $d[I]/dt \approx 0$
  \item a catalyst lowers $E_a$ by a different path; it does not change $K$
\end{itemize}

\knotequestion[\knotetime{38:05}]{Why can a rate law contain the concentration of a species that is not in the overall equation?}
\end{knotechunk}

\begin{knotechunk}
\section{What was said}
\begin{knotetranscript}
\knotetime{14:00:40}Kinetics is about how fast, thermodynamics only tells you whether.

\knotetime{14:04:20}The big warning: you cannot read the orders off the balanced equation, you have to measure them.

\knotetime{14:10:30}The easiest way to find the order is to integrate the rate law and see which plot comes out straight.

\knotetime{14:16:00}For first order the half-life doesn't depend on how much you start with, that's the signature.

\knotetime{14:22:10}Now temperature. Arrhenius noticed that log k against one over T is a straight line.

\knotetime{14:27:40}Take fifty kilojoules per mole and warm by ten degrees, the rate constant almost doubles.

\knotetime{14:34:00}Real reactions happen in steps, and the slowest step controls the overall rate.

\knotetime{14:38:20}A catalyst gives a lower barrier, but the equilibrium constant stays exactly the same.
\end{knotetranscript}
\end{knotechunk}
\end{document}
